\documentclass[a4paper,11pt]{article}

% ---------- Page and typography ----------
\usepackage[a4paper,left=1.25cm,right=1.25cm,top=1.65cm,bottom=1.75cm]{geometry}
\usepackage[T1]{fontenc}
\usepackage[utf8]{inputenc}
\usepackage{lmodern}
\usepackage{microtype}
\usepackage{graphicx}
\usepackage{hyperref}
\usepackage{xurl}
\usepackage{longtable}
\usepackage{booktabs}
\usepackage{array}
\usepackage{calc}
\usepackage{enumitem}
\usepackage{caption}
\usepackage{amssymb}
\usepackage{ragged2e}
\usepackage{textcomp}
\usepackage{adjustbox}
\renewcommand{\_}{\textunderscore\allowbreak}

% ---------- General layout ----------
\setlength{\parindent}{0pt}
\setlength{\parskip}{0.55em}
\setlength{\emergencystretch}{5em}
\sloppy
\justifying

% Lists: compact but readable
\setlist[itemize]{leftmargin=1.1cm,itemsep=0.12em,topsep=0.2em}
\setlist[enumerate]{leftmargin=1.1cm,itemsep=0.12em,topsep=0.2em}

% Tables: clean, centered, readable. No random landscape pages.
\setlength{\tabcolsep}{3.2pt}
\renewcommand{\arraystretch}{1.25}
\captionsetup{justification=centering,singlelinecheck=false}

% Pandoc compatibility
\providecommand{\tightlist}{%
  \setlength{\itemsep}{0pt}\setlength{\parskip}{0pt}}

% Center all longtables and keep readable spacing. Wide tables use smaller text,
% but are not squeezed into unreadable tiny text.
\let\oldlongtable\longtable
\let\endoldlongtable\endlongtable
\renewenvironment{longtable}{%
  \begingroup\centering\scriptsize
  \setlength{\tabcolsep}{2.4pt}%
  \renewcommand{\arraystretch}{1.28}%
  \oldlongtable%
}{%
  \endoldlongtable\endgroup%
}

% Regression outputs in verbatim blocks: compact but readable
\makeatletter
\def\verbatim{\footnotesize\@verbatim \frenchspacing\@vobeyspaces \@xverbatim}
\makeatother

\hypersetup{
  colorlinks=true,
  linkcolor=black,
  urlcolor=blue,
  citecolor=black
}

\begin{document}

\begin{center}
    {\LARGE\bfseries Econometrics}\par
    \vspace{0.6cm}
    {\large Rodrigue Mieuzet}\par
    {\large Nam Khanh NGUYEN}\par
    \vspace{0.35cm}
    {June 2026}\par
\end{center}
\vspace{0.5cm}

\hypertarget{panel-data-econometrics-homework-partial-replication}{%
\section{Panel Data Econometrics Homework: Partial
Replication}\label{panel-data-econometrics-homework-partial-replication}}

\hypertarget{name-of-student}{%
\subsection{Name of student}\label{name-of-student}}

Rodrigue Mieuzet and Nam Khanh NGUYEN.

\hypertarget{identification-of-the-replicated-paper}{%
\subsection{1. Identification of the replicated
paper}\label{identification-of-the-replicated-paper}}

\textbf{APA reference.} Liebersohn, J., \& Rothstein, J. (2025).
\emph{Household mobility and mortgage rate lock}. \emph{Journal of
Financial Economics}, article 103973.
https://doi.org/10.1016/j.jfineco.2024.103973

\textbf{Link to the paper.}
https://doi.org/10.1016/j.jfineco.2024.103973

\hypertarget{table-1-key-correlation-of-the-paper}{%
\subsubsection{Table 1: Key correlation of the
paper}\label{table-1-key-correlation-of-the-paper}}

\begin{itemize}
\tightlist
\item
  Dependent variable Y: \texttt{move\_flag}, an indicator equal to 1
  when a household changes ZIP code between quarters.
\item
  Name and link to the website of data source for updates for Y:
  synthetic replication package built from UC-CCP-style credit registry
  data, public access through Mendeley:
  https://doi.org/10.17632/74sfv9kx9n.1
\item
  Key preferred explanatory variable X: \texttt{rate\_gap}, computed as
  the current market mortgage rate minus the household's implied
  mortgage rate on the outstanding mortgage.
\item
  Name and link to the website of data source for updates for X: Freddie
  Mac mortgage rates through FRED (\texttt{MORTGAGE30US}):
  https://fred.stlouisfed.org/series/MORTGAGE30US
\end{itemize}

\textbf{Selected level(s).} Households.

\textbf{Link to the data set.} https://doi.org/10.17632/74sfv9kx9n.1

\textbf{Links to the data set sources.}

\begin{itemize}
\tightlist
\item
  Replication package data: https://doi.org/10.17632/74sfv9kx9n.1
\item
  FRED mortgage rate series:
  https://fred.stlouisfed.org/series/MORTGAGE30US
\item
  Paper DOI: https://doi.org/10.1016/j.jfineco.2024.103973
\end{itemize}

\hypertarget{abstract-100-words}{%
\subsection{2. Abstract (100 words)}\label{abstract-100-words}}

This partial replication studies whether households move less when their
existing mortgage rate is below the current market rate. Using the
synthetic replication package for Liebersohn and Rothstein, I build a
quarterly household panel from 2012Q1 to 2024Q3, merge mobility events
with mortgage characteristics, and construct a household-level rate-gap
measure. I document panel structure, variable transformations,
descriptive statistics, heterogeneity, and standard panel estimators.
Across between, within, two-way fixed effects, and first-difference
regressions, the estimated relationship between mobility and the rate
gap is small and generally statistically insignificant in this synthetic
sample. The exercise remains informative for panel-data workflow and
diagnostics.

\hypertarget{introduction}{%
\subsection{3. Introduction}\label{introduction}}

The main contribution of this partial replication is methodological
rather than a full reproduction of the authors' confidential-data hazard
framework. The public package contains synthetic data with the same
structure as the original registry-based data, so the present exercise
focuses on building a coherent household-quarter panel, implementing
panel transformations, and checking what survives under between, within,
first-difference, and two-way fixed-effects approaches. This is already
useful because the econometric template asks for an ordered diagnostic
study of the panel before estimation.

Relative to the paper, the present work highlights three things not
emphasized in the published result. First, the synthetic panel is almost
ideal from a teaching perspective: after selection it is balanced, has
no holes, and therefore allows clean comparisons between within, FD and
TWFE transformations. Second, nearly all of the variance of the key
variables comes from the within dimension, which means the key
identifying variation is essentially time variation within households
rather than persistent cross-sectional heterogeneity. Third, once the
data are reduced to this simplified binary mobility outcome and linear
panel estimators, the estimated rate-gap effect is weak.

This does not invalidate the paper. Instead, it clarifies that the
published paper relies on a richer hazard-based design, broader
controls, and confidential micro-data structure that cannot be perfectly
mirrored here. The value of the present replication is therefore to show
a disciplined panel-data workflow, to quantify the differences between
transformed variables, and to explain carefully where the simplified
public exercise is aligned with the paper and where it is not.

\hypertarget{table-2-key-items-found-by-order-of-interest-of-this-partial-replication}{%
\subsubsection{Table 2: Key items found by order of interest of this
partial
replication}\label{table-2-key-items-found-by-order-of-interest-of-this-partial-replication}}

\begin{enumerate}
\def\labelenumi{\arabic{enumi}.}
\tightlist
\item
  The selected sample is fully balanced after dynamic-panel selection:
  1,000 households observed quarterly from 2012Q1 to 2024Q3.
\item
  The within variance share is extremely high for \texttt{move\_flag},
  \texttt{rate\_gap}, \texttt{mortgage\_rate\_annual}, and
  \texttt{loan\_age\_years}, indicating that household time variation
  dominates cross-sectional variation.
\item
  The association between mobility and the rate gap remains weak under
  transformed specifications; it is near zero in within and TWFE and
  mildly negative only in first differences.
\item
  Heterogeneity is substantial across households: some households have
  strongly positive local correlations and others strongly negative
  ones, even when the pooled average is near zero.
\item
  The dynamic FD specification shows limited persistence and weak
  instruments in the simplified Anderson-Hsiao style setup, which
  cautions against over-interpreting dynamic coefficients in this public
  sample.
\end{enumerate}

\hypertarget{database-update}{%
\subsection{4. Database update}\label{database-update}}

The replication package includes a local \texttt{MORTGAGE30US.csv} file
ending on 2023-06-08. I updated the market-rate series using the
official FRED download endpoint before running the final homework
script. This update matters because the mortgage-payment panel extends
to 2024Q3, and without the refreshed FRED series the \texttt{rate\_gap}
variable would be missing near the end of the sample. No new households
were added because the synthetic micro-data are fixed by the replication
package. The main limitation is that the public package still uses
synthetic registry data rather than the confidential underlying data
used in the paper.

\hypertarget{table-3-data-characteristics}{%
\subsubsection{Table 3: Data
characteristics}\label{table-3-data-characteristics}}

\begin{center}
\begingroup
\small
\setlength{\tabcolsep}{4pt}
\renewcommand{\arraystretch}{1.22}
\begin{adjustbox}{max width=\textwidth}
\begin{tabular}{@{}
  >{\raggedright\arraybackslash}p{(\columnwidth - 2\tabcolsep) * \real{0.50}}
  >{\raggedright\arraybackslash}p{(\columnwidth - 2\tabcolsep) * \real{0.50}}@{}}
\toprule
Item & Value \\
\midrule
Maximal time span & 2012Q1 to 2024Q3 \\
Frequency & Quarterly \\
Number of individuals & Ntotal = 1000 \\
Maximal number of time observations for an individual & Tmax = 51 \\
Balanced or unbalanced panel & Balanced \\
Update & FRED mortgage rates updated beyond the local 2023-06-08 cutoff
so that market-rate coverage matches the panel through 2024Q3 \\
\bottomrule
\end{tabular}
\end{adjustbox}
\endgroup
\end{center}

\hypertarget{panel-data-sample-selection}{%
\subsection{5. Panel data sample
selection}\label{panel-data-sample-selection}}

The selection rule is at least three consecutive observations of the
dependent variable for each household. In this synthetic mortgage panel,
every household satisfies this rule.

\textbf{Panel data sample selection.}

\begin{itemize}
\tightlist
\item
  Individuals excluded from the panel data selection: 0
\item
  Individuals with at least three consecutive observations for the
  dependent variable: N1 = 1000
\item
  Ntotal - N1 = 0
\item
  Maximal time span for an individual: 2012Q1 to 2024Q3
\item
  List of excluded individuals: None
\item
  List of kept individuals: omitted here because there are 1,000
  households; the full list is the set of all \texttt{personid} in
  \texttt{panel\_data.csv}.
\end{itemize}

\textbf{Comment on sample-selection bias.} There is no observed
selection bias induced by the three-consecutive-observations rule in
this synthetic sample because no household is dropped. This is
convenient for pedagogy but also a limitation: the public synthetic data
understate the attrition and irregular observation issues commonly found
in real household panels.

\hypertarget{sample-selection-within-an-unbalanced-panel}{%
\subsection{6. Sample selection within an unbalanced
panel}\label{sample-selection-within-an-unbalanced-panel}}

Although the template asks for the analysis of an unbalanced panel, the
selected sample is in fact balanced. Therefore the diagnostics below are
still reported, but they mechanically show no attrition and no holes.

\hypertarget{number-of-individuals-per-date}{%
\subsubsection{Number of individuals per
date}\label{number-of-individuals-per-date}}

\begin{center}
\begingroup
\small
\setlength{\tabcolsep}{5pt}
\renewcommand{\arraystretch}{1.25}
\begin{adjustbox}{max width=\textwidth}
\begin{tabular}{lrrrrrrrrrrrrr}
\toprule
Year & 2012 & 2013 & 2014 & 2015 & 2016 & 2017 & 2018 & 2019 & 2020 & 2021 & 2022 & 2023 & 2024 \\
\midrule
Number of individuals & 1000 & 1000 & 1000 & 1000 & 1000 & 1000 & 1000 & 1000 & 1000 & 1000 & 1000 & 1000 & 1000 \\
\bottomrule
\end{tabular}
\end{adjustbox}
\endgroup
\end{center}

\hypertarget{number-of-individuals-with-the-same-number-of-temporal-observations}{%
\subsubsection{Number of individuals with the same number of temporal
observations}\label{number-of-individuals-with-the-same-number-of-temporal-observations}}

\begin{center}
\begingroup
\footnotesize
\setlength{\tabcolsep}{4pt}
\renewcommand{\arraystretch}{1.22}
\begin{adjustbox}{max width=\textwidth}
\begin{tabular}{@{}lr@{}}
\toprule
quarter & number\_of\_individuals \\
\midrule
51 & 1000 \\
\bottomrule
\end{tabular}
\end{adjustbox}
\endgroup
\end{center}

\begin{itemize}
\tightlist
\item
  Number of individuals with holes: 0
\item
  Proportion with holes: 0.00\%
\item
  List with holes: {[}{]}
\end{itemize}

\textbf{Comment.} There is no attrition at the beginning or at the end
of the sample and no discontinuity pattern. This again reflects the
synthetic design of the teaching data rather than the full complexity of
the confidential panel behind the paper.

\hypertarget{panel-data-variables-classification}{%
\subsection{7. Panel data variables
classification}\label{panel-data-variables-classification}}

\hypertarget{table-4-variable-count-by-categories}{%
\subsubsection{Table 4: Variable count by
categories}\label{table-4-variable-count-by-categories}}

\begin{itemize}
\tightlist
\item
  Variables varying with two indices: rate\_gap, mortgage\_rate\_annual,
  move\_flag, log\_balance, balance\_to\_orig, loan\_age\_years
\item
  Time-invariant variables (100\% between variance): None in the
  selected variable set
\item
  Individual-invariant variables (0\% between variance): market\_rate,
  trend
\item
  Number K = 6
\item
  K1 = 0
\item
  K2 = 2
\end{itemize}

\hypertarget{table-5-time-varying-variables-sorted-by-within-share}{%
\subsubsection{Table 5: Time-varying variables sorted by within
share}\label{table-5-time-varying-variables-sorted-by-within-share}}

\begin{center}
\begingroup
\scriptsize
\setlength{\tabcolsep}{3pt}
\renewcommand{\arraystretch}{1.22}
\begin{adjustbox}{max width=\textwidth}
\begin{tabular}{@{}lrrrrrrr@{}}
\toprule
variable & N & NT & NT\_over\_N & overall\_var & between\_var &
within\_var & within\_share \\
\midrule
rate\_gap & 1000 & 48737 & 48.737000 & 1.889807 & 0.011398 & 1.919089 &
1.015495 \\
mortgage\_rate\_annual & 1000 & 48737 & 48.737000 & 0.456431 & 0.010072
& 0.457020 & 1.001292 \\
move\_flag & 1000 & 51000 & 51.000000 & 0.125795 & 0.002460 & 0.125802 &
1.000055 \\
log\_balance & 1000 & 51000 & 51.000000 & 2.032352 & 0.045471 & 2.026625
& 0.997182 \\
balance\_to\_orig & 1000 & 51000 & 51.000000 & 44131.307643 &
2576.719555 & 42387.425478 & 0.960484 \\
loan\_age\_years & 1000 & 51000 & 51.000000 & 2.393774 & 0.510415 &
1.921499 & 0.802707 \\
\bottomrule
\end{tabular}
\end{adjustbox}
\endgroup
\end{center}


\hypertarget{time-invariant-variables-100-between-variance}{%
\subsubsection{Time-invariant variables (100\% between
variance)}\label{time-invariant-variables-100-between-variance}}

None in the selected set.

\hypertarget{individual-invariant-variables-0-between-variance}{%
\subsubsection{Individual-invariant variables (0\% between
variance)}\label{individual-invariant-variables-0-between-variance}}

\begin{center}
\begingroup
\scriptsize
\setlength{\tabcolsep}{3pt}
\renewcommand{\arraystretch}{1.22}
\begin{adjustbox}{max width=\textwidth}
\begin{tabular}{@{}lrrrrrrr@{}}
\toprule
variable & N & NT & NT\_over\_N & overall\_var & between\_var &
within\_var & within\_share \\
\midrule
market\_rate & 1000 & 51000 & 51.000000 & 1.427681 & 0.000000 & 1.456206
& 1.019980 \\
trend & 1000 & 51000 & 51.000000 & 216.670915 & 0.000000 & 221.000000 &
1.019980 \\
\bottomrule
\end{tabular}
\end{adjustbox}
\endgroup
\end{center}


\textbf{Comment.} The key dependent variable \texttt{move\_flag} and the
key explanatory variable \texttt{rate\_gap} both have very high within
shares. This means that the informative variation for the homework
mostly comes from the evolution of a household over time rather than
from persistent differences across households. The common mortgage-rate
series (\texttt{market\_rate}) and the deterministic time trend are
individual-invariant variables, so they are wiped out by time dummies
and/or by TWFE transformation.

\hypertarget{between-plus-one-way-within-decomposition}{%
\subsection{8. Between plus one-way within
decomposition}\label{between-plus-one-way-within-decomposition}}

\hypertarget{descriptive-statistics-for-the-dependent-variable-move_flag}{%
\subsubsection{\texorpdfstring{Descriptive statistics for the dependent
variable
\texttt{move\_flag}}{Descriptive statistics for the dependent variable move\_flag}}\label{descriptive-statistics-for-the-dependent-variable-move_flag}}

\begin{center}
\begingroup
\scriptsize
\setlength{\tabcolsep}{3pt}
\renewcommand{\arraystretch}{1.22}
\begin{adjustbox}{max width=\textwidth}
\begin{tabular}{@{}lrrrrrrrrrrrrr@{}}
\toprule
transform & n & mean & median & sd & se & min & q1 & q3 & max & skew &
kurtosis & std\_min & std\_max \\
\midrule
between & 51000 & 0.147569 & 0.137255 & 0.049571 & 0.049571 & 0.019608 &
0.117647 & 0.176471 & 0.333333 & 0.387867 & 3.335060 & -2.581353 &
3.747432 \\
within & 51000 & -0.000000 & -0.137255 & 0.351194 & 0.351194 & -0.333333
& -0.176471 & -0.098039 & 0.980392 & 1.929329 & 4.877890 & -0.949143 &
2.791596 \\
fd & 50000 & 0.003060 & 0.000000 & 0.503007 & 0.503007 & -1.000000 &
0.000000 & 0.000000 & 1.000000 & 0.005794 & 3.952401 & -1.994129 &
1.981962 \\
twfe & 51000 & 0.000000 & -0.128686 & 0.346629 & 0.346629 & -0.614549 &
-0.171902 & -0.074863 & 1.070353 & 1.874439 & 4.867251 & -1.772930 &
3.087892 \\
\bottomrule
\end{tabular}
\end{adjustbox}
\endgroup
\end{center}


\hypertarget{descriptive-statistics-for-the-preferred-explanatory-variable-rate_gap}{%
\subsubsection{\texorpdfstring{Descriptive statistics for the preferred
explanatory variable
\texttt{rate\_gap}}{Descriptive statistics for the preferred explanatory variable rate\_gap}}\label{descriptive-statistics-for-the-preferred-explanatory-variable-rate_gap}}

\begin{center}
\begingroup
\scriptsize
\setlength{\tabcolsep}{3pt}
\renewcommand{\arraystretch}{1.22}
\begin{adjustbox}{max width=\textwidth}
\begin{tabular}{@{}lrrrrrrrrrrrrr@{}}
\toprule
transform & n & mean & median & sd & se & min & q1 & q3 & max & skew &
kurtosis & std\_min & std\_max \\
\midrule
between & 51000 & 4.092290 & 4.114000 & 0.106710 & 0.106710 & 3.576468 &
4.037769 & 4.168272 & 4.290295 & -1.112028 & 4.861055 & -4.833877 &
1.855547 \\
within & 48737 & 0.000000 & -0.282254 & 1.370572 & 1.370572 & -11.919270
& -0.686838 & 0.303433 & 3.655761 & -0.226394 & 8.734244 & -8.696563 &
2.667325 \\
fd & 45668 & 0.049493 & -0.003593 & 1.036280 & 1.036280 & -11.849085 &
-0.242456 & 0.290738 & 12.431947 & -0.034364 & 40.936045 & -11.482011 &
11.948947 \\
twfe & 48737 & -0.000000 & 0.110874 & 0.668686 & 0.668686 & -11.335833 &
0.015642 & 0.199821 & 0.795920 & -8.586655 & 100.252310 & -16.952390 &
1.190274 \\
\bottomrule
\end{tabular}
\end{adjustbox}
\endgroup
\end{center}


The within-transformed mobility indicator is highly non-normal and
discrete, with many observations concentrated near household means. The
between distribution of \texttt{rate\_gap} is tighter than the within
distribution, which is consistent with the idea that market rates and
loan-specific conditions move substantially over time within households.
The graphs saved in
\texttt{analysis\_outputs/move\_flag\_transforms.png} and
\texttt{analysis\_outputs/rate\_gap\_transforms.png} display the
histogram, KDE, and matching normal approximation requested in the
template.

\hypertarget{first-differences-versus-two-way-fixed-effects}{%
\subsection{9. First differences versus two-way fixed
effects}\label{first-differences-versus-two-way-fixed-effects}}

Because the selected sample is balanced, the closed-form TWFE
transformation \texttt{x\_it\ -\ x\_i.\ -\ x\_.t\ +\ x\_..} is well
defined and coincides with the regression-residual interpretation. For
first differences, the first observation of each individual is missing
by construction. The file \texttt{analysis\_outputs/fd\_first30.csv}
reports the first 30 stacked observations of first differences and their
lags, and one can verify that each individual switch generates missing
first differences rather than cross-household subtraction.

\hypertarget{between-vs-within-correlation-matrix}{%
\subsubsection{Between vs within correlation
matrix}\label{between-vs-within-correlation-matrix}}

\begin{center}
\begingroup
\footnotesize
\setlength{\tabcolsep}{4pt}
\renewcommand{\arraystretch}{1.22}
\begin{adjustbox}{max width=\textwidth}
\begin{tabular}{@{}lrrrrr@{}}
\toprule
index & move\_flag\_between & rate\_gap\_between &
balance\_to\_orig\_between & loan\_age\_years\_between &
log\_balance\_between \\
\midrule
move\_flag\_between & 1.000000 & 0.052314 & 0.004630 & -0.416690 &
0.001786 \\
rate\_gap\_between & 0.052314 & 1.000000 & -0.003660 & -0.028802 &
0.473497 \\
balance\_to\_orig\_between & 0.004630 & -0.003660 & 1.000000 & -0.001591
& -0.083276 \\
loan\_age\_years\_between & -0.416690 & -0.028802 & -0.001591 & 1.000000
& 0.024215 \\
log\_balance\_between & 0.001786 & 0.473497 & -0.083276 & 0.024215 &
1.000000 \\
\bottomrule
\end{tabular}
\end{adjustbox}
\endgroup
\end{center}

\begin{center}
\begingroup
\footnotesize
\setlength{\tabcolsep}{4pt}
\renewcommand{\arraystretch}{1.22}
\begin{adjustbox}{max width=\textwidth}
\begin{tabular}{@{}lrrrrrr@{}}
\toprule
index & move\_flag\_within & rate\_gap\_within &
balance\_to\_orig\_within & loan\_age\_years\_within &
log\_balance\_within & trend\_within \\
\midrule
move\_flag\_within & 1.000000 & 0.004038 & -0.007428 & -0.174514 &
0.004541 & 0.014490 \\
rate\_gap\_within & 0.004038 & 1.000000 & -0.002453 & 0.057984 &
0.243703 & 0.466081 \\
balance\_to\_orig\_within & -0.007428 & -0.002453 & 1.000000 & -0.005025
& 0.022447 & -0.004764 \\
loan\_age\_years\_within & -0.174514 & 0.057984 & -0.005025 & 1.000000 &
-0.005559 & 0.252264 \\
log\_balance\_within & 0.004541 & 0.243703 & 0.022447 & -0.005559 &
1.000000 & 0.000096 \\
trend\_within & 0.014490 & 0.466081 & -0.004764 & 0.252264 & 0.000096 &
1.000000 \\
\bottomrule
\end{tabular}
\end{adjustbox}
\endgroup
\end{center}


\hypertarget{twfe-vs-fd-correlation-matrix}{%
\subsubsection{TWFE vs FD correlation
matrix}\label{twfe-vs-fd-correlation-matrix}}

\begin{center}
\begingroup
\footnotesize
\setlength{\tabcolsep}{4pt}
\renewcommand{\arraystretch}{1.22}
\begin{adjustbox}{max width=\textwidth}
\begin{tabular}{@{}lrrrrr@{}}
\toprule
index & move\_flag\_twfe & rate\_gap\_twfe & balance\_to\_orig\_twfe &
loan\_age\_years\_twfe & log\_balance\_twfe \\
\midrule
move\_flag\_twfe & 1.000000 & 0.001892 & -0.007375 & -0.198220 &
0.004439 \\
rate\_gap\_twfe & 0.001892 & 1.000000 & 0.006838 & -0.000247 &
0.482349 \\
balance\_to\_orig\_twfe & -0.007375 & 0.006838 & 1.000000 & -0.003654 &
0.022366 \\
loan\_age\_years\_twfe & -0.198220 & -0.000247 & -0.003654 & 1.000000 &
-0.003955 \\
log\_balance\_twfe & 0.004439 & 0.482349 & 0.022366 & -0.003955 &
1.000000 \\
\bottomrule
\end{tabular}
\end{adjustbox}
\endgroup
\end{center}

\begin{center}
\begingroup
\scriptsize
\setlength{\tabcolsep}{3pt}
\renewcommand{\arraystretch}{1.22}
\begin{adjustbox}{max width=\textwidth}
\begin{tabular}{@{}lrrrrrrr@{}}
\toprule
index & move\_flag\_fd & rate\_gap\_fd & balance\_to\_orig\_fd &
loan\_age\_years\_fd & log\_balance\_fd & move\_flag\_fd\_lag1 &
rate\_gap\_fd\_lag1 \\
\midrule
move\_flag\_fd & 1.000000 & -0.007410 & -0.014307 & -0.305840 & 0.001279
& -0.503768 & 0.000742 \\
rate\_gap\_fd & -0.007410 & 1.000000 & 0.006575 & -0.006187 & 0.450417 &
0.006452 & -0.412148 \\
balance\_to\_orig\_fd & -0.014307 & 0.006575 & 1.000000 & 0.007372 &
0.029813 & 0.002737 & -0.001369 \\
loan\_age\_years\_fd & -0.305840 & -0.006187 & 0.007372 & 1.000000 &
-0.003815 & 0.005156 & -0.007473 \\
log\_balance\_fd & 0.001279 & 0.450417 & 0.029813 & -0.003815 & 1.000000
& 0.004003 & -0.230097 \\
move\_flag\_fd\_lag1 & -0.503768 & 0.006452 & 0.002737 & 0.005156 &
0.004003 & 1.000000 & -0.007216 \\
rate\_gap\_fd\_lag1 & 0.000742 & -0.412148 & -0.001369 & -0.007473 &
-0.230097 & -0.007216 & 1.000000 \\
\bottomrule
\end{tabular}
\end{adjustbox}
\endgroup
\end{center}


The simple correlations between transformed mobility and transformed
rate-gap variables remain weak in absolute value, usually below 0.1,
which is exactly the kind of diagnostic the template asks us to note
before moving to regressions. The strongest correlations are mostly
among explanatory variables and their lags rather than between the
dependent variable and the key explanatory variable.

\hypertarget{investigating-bivariate-heterogeneity-by-individuals-for-twfe-and-fd}{%
\subsection{10. Investigating bivariate heterogeneity by individuals for
TWFE and
FD}\label{investigating-bivariate-heterogeneity-by-individuals-for-twfe-and-fd}}

\hypertarget{fd-heterogeneity-table-top-15-households}{%
\subsubsection{FD heterogeneity table (top 15
households)}\label{fd-heterogeneity-table-top-15-households}}

\begin{center}
\begingroup
\footnotesize
\setlength{\tabcolsep}{4pt}
\renewcommand{\arraystretch}{1.22}
\begin{adjustbox}{max width=\textwidth}
\begin{tabular}{@{}lrrrrrr@{}}
\toprule
personid & T\_i & corr & sd\_y & sd\_x & sd\_ratio & beta \\
\midrule
A865000 & 48 & 0.487346 & 0.291730 & 0.716015 & 0.407436 & 0.198562 \\
A114000 & 46 & 0.457269 & 0.493924 & 2.083628 & 0.237050 & 0.108396 \\
A749000 & 46 & 0.433213 & 0.446673 & 1.092391 & 0.408895 & 0.177139 \\
A633000 & 48 & 0.410655 & 0.545777 & 0.608111 & 0.897495 & 0.368561 \\
A464000 & 41 & 0.408283 & 0.523823 & 0.441959 & 1.185229 & 0.483910 \\
A806000 & 42 & 0.402288 & 0.538851 & 1.353879 & 0.398005 & 0.160113 \\
A597000 & 48 & 0.399628 & 0.635462 & 0.432059 & 1.470775 & 0.587763 \\
A198000 & 48 & 0.389519 & 0.601048 & 0.898755 & 0.668757 & 0.260494 \\
A963000 & 44 & 0.388515 & 0.681994 & 0.481448 & 1.416547 & 0.550350 \\
A106000 & 46 & 0.378778 & 0.446673 & 1.475461 & 0.302735 & 0.114669 \\
A160000 & 48 & 0.364132 & 0.618853 & 0.564793 & 1.095716 & 0.398986 \\
A417000 & 40 & 0.354812 & 0.552384 & 0.537491 & 1.027710 & 0.364644 \\
A988000 & 48 & 0.352375 & 0.525502 & 0.518943 & 1.012639 & 0.356828 \\
A228000 & 45 & 0.332781 & 0.424026 & 1.820645 & 0.232899 & 0.077504 \\
A577000 & 46 & 0.331839 & 0.596285 & 0.414725 & 1.437784 & 0.477113 \\
\bottomrule
\end{tabular}
\end{adjustbox}
\endgroup
\end{center}


\hypertarget{twfe-heterogeneity-table-top-15-households}{%
\subsubsection{TWFE heterogeneity table (top 15
households)}\label{twfe-heterogeneity-table-top-15-households}}

\begin{center}
\begingroup
\footnotesize
\setlength{\tabcolsep}{4pt}
\renewcommand{\arraystretch}{1.22}
\begin{adjustbox}{max width=\textwidth}
\begin{tabular}{@{}lrrrrrr@{}}
\toprule
personid & T\_i & corr & sd\_y & sd\_x & sd\_ratio & beta \\
\midrule
A963000 & 48 & 0.270483 & 0.458631 & 0.214467 & 2.138473 & 0.578420 \\
A803000 & 49 & 0.264065 & 0.316578 & 0.133697 & 2.367875 & 0.625274 \\
A939000 & 50 & 0.255064 & 0.374256 & 0.251831 & 1.486139 & 0.379060 \\
A632000 & 48 & 0.231929 & 0.413671 & 0.229639 & 1.801395 & 0.417795 \\
A859000 & 48 & 0.226349 & 0.381332 & 0.190292 & 2.003933 & 0.453587 \\
A633000 & 50 & 0.219752 & 0.353331 & 0.336027 & 1.051497 & 0.231069 \\
A266000 & 49 & 0.218839 & 0.333143 & 0.104639 & 3.183735 & 0.696725 \\
A401000 & 51 & 0.216054 & 0.376469 & 0.401410 & 0.937867 & 0.202630 \\
A931000 & 49 & 0.215674 & 0.333057 & 0.402531 & 0.827409 & 0.178450 \\
A417000 & 46 & 0.214949 & 0.420469 & 0.280954 & 1.496573 & 0.321686 \\
A250000 & 47 & 0.207706 & 0.384815 & 0.201274 & 1.911899 & 0.397112 \\
A820000 & 50 & 0.206220 & 0.389418 & 0.180982 & 2.151690 & 0.443722 \\
A239000 & 47 & 0.205037 & 0.315982 & 0.687771 & 0.459430 & 0.094200 \\
A024000 & 50 & 0.203571 & 0.406944 & 0.280416 & 1.451215 & 0.295426 \\
A083000 & 48 & 0.202623 & 0.311964 & 0.110329 & 2.827565 & 0.572931 \\
\bottomrule
\end{tabular}
\end{adjustbox}
\endgroup
\end{center}


\begin{itemize}
\tightlist
\item
  FD diagnosis: positive correlation households = 288, negative
  correlation households = 243, weak-correlation households = 469.
\item
  TWFE diagnosis: positive correlation households = 346, negative
  correlation households = 196, weak-correlation households = 458.
\end{itemize}

The pooled effect is small partly because heterogeneity is large: some
households show positive local comovement, others negative comovement.
This is precisely why the template asks for household-level
heterogeneity tables rather than relying only on pooled coefficients.
Rolling-window plots for representative households are saved in
\texttt{analysis\_outputs/fd\_rolling\_corr.png} and
\texttt{analysis\_outputs/twfe\_rolling\_corr.png}.

\hypertarget{panel-data-estimates}{%
\subsection{11. Panel data estimates}\label{panel-data-estimates}}

\begin{center}
\begingroup
\footnotesize
\setlength{\tabcolsep}{4pt}
\renewcommand{\arraystretch}{1.22}
\begin{adjustbox}{max width=\textwidth}
\begin{tabular}{@{}lrrrr@{}}
\toprule
model & n\_obs & coef\_rate\_gap & se\_rate\_gap & r\_squared \\
\midrule
between & 1000 & 0.020735 & 0.015261 & 0.175337 \\
within & 48737 & 0.003618 & 0.001194 & 0.030725 \\
mundlak & 48737 & 0.003629 & 0.001209 & 0.033640 \\
twfe & 48737 & 0.000032 & 0.002686 & 0.039371 \\
fd & 45668 & -0.005959 & 0.002360 & 0.092429 \\
\bottomrule
\end{tabular}
\end{adjustbox}
\endgroup
\end{center}

Across between, within, Mundlak, TWFE and first-difference estimators,
the coefficient on \texttt{rate\_gap} remains economically small in the
synthetic sample. The between, within and Mundlak coefficients are
slightly positive, the TWFE coefficient is essentially zero, and only
the first-difference coefficient is negative. This is an important
replication conclusion: the simplified public panel exercise does not
recover a stable reduced-form sign comparable to the full
confidential-data hazard specification.

\hypertarget{dynamic-panel-anderson-hsiao-style-extension}{%
\subsubsection{Dynamic panel / Anderson-Hsiao style
extension}\label{dynamic-panel-anderson-hsiao-style-extension}}

\hypertarget{univariate-statistics-for-dynamic-variables}{%
\paragraph{Univariate statistics for dynamic
variables}\label{univariate-statistics-for-dynamic-variables}}

\begin{center}
\begingroup
\scriptsize
\setlength{\tabcolsep}{3pt}
\renewcommand{\arraystretch}{1.22}
\begin{adjustbox}{max width=\textwidth}
\begin{tabular}{@{}lrrrrrrrrrrrrr@{}}
\toprule
variable & n & mean & median & sd & se & min & q1 & q3 & max & skew &
kurtosis & std\_min & std\_max \\
\midrule
dy & 42799 & 0.001612 & 0.000000 & 0.509109 & 0.509109 & -1.000000 &
0.000000 & 0.000000 & 1.000000 & 0.002718 & 3.858422 & -1.967383 &
1.961050 \\
dy\_l1 & 42799 & 0.001916 & 0.000000 & 0.506093 & 0.506093 & -1.000000 &
0.000000 & 0.000000 & 1.000000 & 0.003424 & 3.904526 & -1.979708 &
1.972137 \\
dx & 42799 & 0.056031 & 0.005064 & 1.043341 & 1.043341 & -11.849085 &
-0.242583 & 0.297063 & 12.431947 & 0.065063 & 40.629890 & -11.410566 &
11.861809 \\
dx\_l1 & 42799 & 0.058741 & 0.007574 & 1.036956 & 1.036956 & -11.849085
& -0.229883 & 0.297836 & 12.431947 & -0.086836 & 41.130615 & -11.483440
& 11.932235 \\
y\_l2 & 42799 & 0.149419 & 0.000000 & 0.356506 & 0.356506 & 0.000000 &
0.000000 & 0.000000 & 1.000000 & 1.966853 & 4.868598 & -0.419122 &
2.385883 \\
x\_l2 & 42799 & 3.991430 & 3.791079 & 1.298371 & 1.298371 & -8.181924 &
3.406824 & 4.286574 & 7.303224 & -0.404210 & 10.995020 & -9.375866 &
2.550730 \\
\bottomrule
\end{tabular}
\end{adjustbox}
\endgroup
\end{center}


\hypertarget{correlation-matrix-for-dynamic-variables-and-instruments}{%
\paragraph{Correlation matrix for dynamic variables and
instruments}\label{correlation-matrix-for-dynamic-variables-and-instruments}}

\begin{center}
\begingroup
\footnotesize
\setlength{\tabcolsep}{4pt}
\renewcommand{\arraystretch}{1.22}
\begin{adjustbox}{max width=\textwidth}
\begin{tabular}{@{}lrrrrrr@{}}
\toprule
index & dy & dy\_l1 & dx & dx\_l1 & y\_l2 & x\_l2 \\
\midrule
dy & 1.000000 & -0.503768 & -0.007410 & 0.000742 & 0.004208 &
-0.004623 \\
dy\_l1 & -0.503768 & 1.000000 & 0.006452 & -0.007216 & -0.706084 &
-0.003264 \\
dx & -0.007410 & 0.006452 & 1.000000 & -0.412148 & 0.004068 &
-0.037537 \\
dx\_l1 & 0.000742 & -0.007216 & -0.412148 & 1.000000 & 0.015926 &
-0.336558 \\
y\_l2 & 0.004208 & -0.706084 & 0.004068 & 0.015926 & 1.000000 &
0.002562 \\
x\_l2 & -0.004623 & -0.003264 & -0.037537 & -0.336558 & 0.002562 &
1.000000 \\
\bottomrule
\end{tabular}
\end{adjustbox}
\endgroup
\end{center}


\begin{itemize}
\tightlist
\item
  ADF test on $\Delta Y$: statistic = -71.831, p-value = 0.0000
\item
  ADF test on $\Delta X$: statistic = -26.107, p-value = 0.0000
\end{itemize}

These pooled ADF tests strongly suggest stationarity of the differenced
series, which is the expected result after first differencing. The
template asked for a panel unit-root test of our choice; in this local
Python environment, the pooled ADF test is the feasible choice.

\hypertarget{ols-dynamic-fd-regression}{%
\paragraph{OLS dynamic FD regression}\label{ols-dynamic-fd-regression}}

\begin{verbatim}
                            OLS Regression Results                            
==============================================================================
Dep. Variable:                     dy   R-squared:                       0.364
Model:                            OLS   Adj. R-squared:                  0.363
Method:                 Least Squares   F-statistic:                     303.6
Date:                Thu, 28 May 2026   Prob (F-statistic):               0.00
Time:                        17:15:49   Log-Likelihood:                -22159.
No. Observations:               42799   AIC:                         4.443e+04
Df Residuals:                   42744   BIC:                         4.491e+04
Df Model:                          54                                         
Covariance Type:                  HC1                                         
=====================================================================================
                        coef    std err          z      P>|z|      [0.025      0.975]
-------------------------------------------------------------------------------------
const                 0.0666      0.007      9.777      0.000       0.053       0.080
dx                   -0.0012      0.003     -0.475      0.635      -0.006       0.004
d_balance_to_orig -2.153e-05   7.62e-06     -2.827      0.005   -3.65e-05    -6.6e-06
d_loan_age_years     -0.2001      0.005    -43.880      0.000      -0.209      -0.191
d_log_balance         0.0009      0.001      0.755      0.450      -0.001       0.003
q_2012Q4              0.0118      0.011      1.078      0.281      -0.010       0.033
q_2013Q1              0.3570      0.016     22.341      0.000       0.326       0.388
q_2013Q2             -0.1538      0.017     -9.202      0.000      -0.187      -0.121
q_2013Q3             -0.2152      0.017    -12.513      0.000      -0.249      -0.181
q_2013Q4             -0.0598      0.016     -3.813      0.000      -0.090      -0.029
q_2014Q1             -0.0584      0.015     -3.836      0.000      -0.088      -0.029
q_2014Q2             -0.0513      0.015     -3.324      0.001      -0.082      -0.021
q_2014Q3             -0.0390      0.015     -2.620      0.009      -0.068      -0.010
q_2014Q4             -0.0405      0.015     -2.629      0.009      -0.071      -0.010
q_2015Q1             -0.0599      0.016     -3.824      0.000      -0.091      -0.029
q_2015Q2             -0.0661      0.016     -4.242      0.000      -0.097      -0.036
q_2015Q3             -0.0399      0.016     -2.516      0.012      -0.071      -0.009
q_2015Q4             -0.0718      0.015     -4.705      0.000      -0.102      -0.042
q_2016Q1             -0.0856      0.016     -5.501      0.000      -0.116      -0.055
q_2016Q2             -0.0692      0.015     -4.478      0.000      -0.099      -0.039
q_2016Q3             -0.0667      0.015     -4.433      0.000      -0.096      -0.037
q_2016Q4             -0.0633      0.015     -4.204      0.000      -0.093      -0.034
q_2017Q1             -0.0462      0.016     -2.982      0.003      -0.077      -0.016
q_2017Q2             -0.0521      0.016     -3.244      0.001      -0.084      -0.021
q_2017Q3             -0.0744      0.016     -4.757      0.000      -0.105      -0.044
q_2017Q4             -0.0686      0.015     -4.442      0.000      -0.099      -0.038
q_2018Q1             -0.0620      0.015     -4.130      0.000      -0.091      -0.033
q_2018Q2             -0.0815      0.016     -5.164      0.000      -0.112      -0.051
q_2018Q3             -0.0957      0.015     -6.273      0.000      -0.126      -0.066
q_2018Q4             -0.0455      0.016     -2.897      0.004      -0.076      -0.015
q_2019Q1             -0.0469      0.016     -2.984      0.003      -0.078      -0.016
q_2019Q2             -0.0648      0.015     -4.239      0.000      -0.095      -0.035
q_2019Q3             -0.0764      0.015     -5.056      0.000      -0.106      -0.047
q_2019Q4             -0.0628      0.015     -4.237      0.000      -0.092      -0.034
q_2020Q1             -0.0803      0.015     -5.502      0.000      -0.109      -0.052
q_2020Q2             -0.0523      0.015     -3.469      0.001      -0.082      -0.023
q_2020Q3             -0.0486      0.016     -3.097      0.002      -0.079      -0.018
q_2020Q4             -0.0783      0.015     -5.056      0.000      -0.109      -0.048
q_2021Q1             -0.0497      0.016     -3.042      0.002      -0.082      -0.018
q_2021Q2             -0.0748      0.015     -4.844      0.000      -0.105      -0.045
q_2021Q3             -0.0811      0.015     -5.351      0.000      -0.111      -0.051
q_2021Q4             -0.0623      0.016     -4.017      0.000      -0.093      -0.032
q_2022Q1             -0.0631      0.016     -3.997      0.000      -0.094      -0.032
q_2022Q2             -0.0427      0.017     -2.547      0.011      -0.076      -0.010
q_2022Q3             -0.0642      0.016     -3.985      0.000      -0.096      -0.033
q_2022Q4             -0.0915      0.016     -5.612      0.000      -0.123      -0.060
q_2023Q1             -0.0833      0.016     -5.158      0.000      -0.115      -0.052
q_2023Q2             -0.0725      0.015     -4.883      0.000      -0.102      -0.043
q_2023Q3             -0.0728      0.015     -4.780      0.000      -0.103      -0.043
q_2023Q4             -0.0503      0.015     -3.300      0.001      -0.080      -0.020
q_2024Q1             -0.0542      0.015     -3.552      0.000      -0.084      -0.024
q_2024Q2             -0.0691      0.016     -4.445      0.000      -0.100      -0.039
q_2024Q3             -0.0637      0.016     -3.999      0.000      -0.095      -0.032
dy_l1                -0.5075      0.005   -112.051      0.000      -0.516      -0.499
dx_l1                 0.0004      0.002      0.187      0.851      -0.004       0.005
==============================================================================
Omnibus:                     2923.603   Durbin-Watson:                   2.272
Prob(Omnibus):                  0.000   Jarque-Bera (JB):             5065.037
Skew:                           0.520   Prob(JB):                         0.00
Kurtosis:                       4.326   Cond. No.                     1.29e+04
==============================================================================

Notes:
[1] Standard Errors are heteroscedasticity robust (HC1)
[2] The condition number is large, 1.29e+04. This might indicate that there are
strong multicollinearity or other numerical problems.
\end{verbatim}

\hypertarget{iv-dynamic-fd-regression-levels-instruments-y_t-2-and-x_t-2}{%
\paragraph{\texorpdfstring{IV dynamic FD regression (levels instruments
\texttt{Y\_\{t-2\}} and
\texttt{X\_\{t-2\}})}{IV dynamic FD regression (levels instruments Y\_\{t-2\} and X\_\{t-2\})}}\label{iv-dynamic-fd-regression-levels-instruments-y_t-2-and-x_t-2}}

\begin{verbatim}
                          IV2SLS Regression Results                           
==============================================================================
Dep. Variable:                     dy   R-squared:                       0.143
Model:                         IV2SLS   Adj. R-squared:                  0.142
Method:                     Two Stage   F-statistic:                     106.7
                        Least Squares   Prob (F-statistic):               0.00
Date:                Thu, 28 May 2026                                         
Time:                        17:15:49                                         
No. Observations:               42799                                         
Df Residuals:                   42744                                         
Df Model:                          54                                         
=====================================================================================
                        coef    std err          t      P>|t|      [0.025      0.975]
-------------------------------------------------------------------------------------
const                 0.0533      0.016      3.340      0.001       0.022       0.085
dx                   -0.0031      0.003     -0.982      0.326      -0.009       0.003
d_balance_to_orig -2.449e-05   8.78e-06     -2.788      0.005   -4.17e-05   -7.27e-06
d_loan_age_years     -0.2012      0.003    -67.917      0.000      -0.207      -0.195
d_log_balance         0.0014      0.001      1.061      0.289      -0.001       0.004
q_2012Q4              0.0209      0.023      0.927      0.354      -0.023       0.065
q_2013Q1              0.3560      0.022     15.832      0.000       0.312       0.400
q_2013Q2             -0.3367      0.023    -14.827      0.000      -0.381      -0.292
q_2013Q3             -0.0442      0.023     -1.918      0.055      -0.089       0.001
q_2013Q4             -0.0466      0.023     -2.044      0.041      -0.091      -0.002
q_2014Q1             -0.0325      0.023     -1.441      0.150      -0.077       0.012
q_2014Q2             -0.0395      0.022     -1.757      0.079      -0.084       0.005
q_2014Q3             -0.0248      0.022     -1.103      0.270      -0.069       0.019
q_2014Q4             -0.0331      0.023     -1.469      0.142      -0.077       0.011
q_2015Q1             -0.0509      0.023     -2.259      0.024      -0.095      -0.007
q_2015Q2             -0.0445      0.023     -1.973      0.048      -0.089      -0.000
q_2015Q3             -0.0241      0.023     -1.071      0.284      -0.068       0.020
q_2015Q4             -0.0703      0.022     -3.125      0.002      -0.114      -0.026
q_2016Q1             -0.0626      0.023     -2.780      0.005      -0.107      -0.018
q_2016Q2             -0.0507      0.023     -2.251      0.024      -0.095      -0.007
q_2016Q3             -0.0515      0.022     -2.291      0.022      -0.095      -0.007
q_2016Q4             -0.0467      0.023     -2.071      0.038      -0.091      -0.002
q_2017Q1             -0.0372      0.023     -1.642      0.101      -0.082       0.007
q_2017Q2             -0.0473      0.023     -2.091      0.037      -0.092      -0.003
q_2017Q3             -0.0622      0.023     -2.761      0.006      -0.106      -0.018
q_2017Q4             -0.0494      0.023     -2.187      0.029      -0.094      -0.005
q_2018Q1             -0.0490      0.023     -2.170      0.030      -0.093      -0.005
q_2018Q2             -0.0689      0.023     -3.042      0.002      -0.113      -0.025
q_2018Q3             -0.0806      0.023     -3.549      0.000      -0.125      -0.036
q_2018Q4             -0.0215      0.023     -0.943      0.346      -0.066       0.023
q_2019Q1             -0.0429      0.023     -1.893      0.058      -0.087       0.002
q_2019Q2             -0.0539      0.023     -2.385      0.017      -0.098      -0.010
q_2019Q3             -0.0639      0.023     -2.828      0.005      -0.108      -0.020
q_2019Q4             -0.0400      0.023     -1.772      0.076      -0.084       0.004
q_2020Q1             -0.0639      0.023     -2.829      0.005      -0.108      -0.020
q_2020Q2             -0.0286      0.023     -1.261      0.207      -0.073       0.016
q_2020Q3             -0.0480      0.023     -2.123      0.034      -0.092      -0.004
q_2020Q4             -0.0602      0.023     -2.667      0.008      -0.105      -0.016
q_2021Q1             -0.0372      0.023     -1.651      0.099      -0.081       0.007
q_2021Q2             -0.0724      0.023     -3.207      0.001      -0.117      -0.028
q_2021Q3             -0.0548      0.023     -2.428      0.015      -0.099      -0.011
q_2021Q4             -0.0407      0.023     -1.799      0.072      -0.085       0.004
q_2022Q1             -0.0533      0.023     -2.310      0.021      -0.098      -0.008
q_2022Q2             -0.0281      0.024     -1.182      0.237      -0.075       0.019
q_2022Q3             -0.0688      0.024     -2.910      0.004      -0.115      -0.022
q_2022Q4             -0.0790      0.023     -3.419      0.001      -0.124      -0.034
q_2023Q1             -0.0615      0.023     -2.673      0.008      -0.107      -0.016
q_2023Q2             -0.0591      0.023     -2.613      0.009      -0.103      -0.015
q_2023Q3             -0.0507      0.023     -2.225      0.026      -0.095      -0.006
q_2023Q4             -0.0440      0.023     -1.931      0.054      -0.089       0.001
q_2024Q1             -0.0446      0.023     -1.973      0.048      -0.089      -0.000
q_2024Q2             -0.0631      0.023     -2.790      0.005      -0.107      -0.019
q_2024Q3             -0.0416      0.023     -1.839      0.066      -0.086       0.003
dy_l1                -0.0290      0.006     -4.501      0.000      -0.042      -0.016
dx_l1                 0.0041      0.003      1.213      0.225      -0.003       0.011
==============================================================================
Omnibus:                     1304.124   Durbin-Watson:                   2.863
Prob(Omnibus):                  0.000   Jarque-Bera (JB):             2769.515
Skew:                          -0.193   Prob(JB):                         0.00
Kurtosis:                       4.185   Cond. No.                     1.29e+04
==============================================================================
\end{verbatim}

\hypertarget{first-stage-regressions}{%
\paragraph{First-stage regressions}\label{first-stage-regressions}}

\begin{itemize}
\tightlist
\item
  First stage for \texttt{$\Delta Y$\_(t-1)}: $R^2$ = 0.5117
\item
  First stage for \texttt{$\Delta X$\_(t-1)}: $R^2$ = 0.7811
\end{itemize}

The first-stage $R^2$ values are not especially low in this simplified
setup, so there is no immediate weak-instrument warning from that
diagnostic alone. Even so, the dynamic IV exercise should still be read
as a pedagogical robustness check rather than a definitive causal
estimate, because the public synthetic panel is not the authors' full
hazard dataset.

\hypertarget{impulse-response-and-long-run-coefficient}{%
\paragraph{Impulse response and long-run
coefficient}\label{impulse-response-and-long-run-coefficient}}

\begin{center}
\begingroup
\footnotesize
\setlength{\tabcolsep}{4pt}
\renewcommand{\arraystretch}{1.22}
\begin{adjustbox}{max width=\textwidth}
\begin{tabular}{@{}lr@{}}
\toprule
horizon & response \\
\midrule
1.000000 & -0.001236 \\
2.000000 & 0.001074 \\
3.000000 & -0.000545 \\
4.000000 & 0.000276 \\
\bottomrule
\end{tabular}
\end{adjustbox}
\endgroup
\end{center}

\begin{itemize}
\tightlist
\item
  Long-run coefficient: -0.000524
\end{itemize}

\hypertarget{optional-time-invariant-variable-estimators}{%
\subsection{12. Optional time-invariant-variable
estimators}\label{optional-time-invariant-variable-estimators}}

Skipped. In the selected working set for this homework, I did not keep a
substantively meaningful time-invariant household characteristic that
could justify a Hausman-Taylor exercise. Forcing an artificial
time-invariant regressor would not improve the econometric quality of
the homework. This is preferable to running a mechanical but
uninformative Hausman-Taylor specification.

\hypertarget{own-estimations-not-done-in-the-paper}{%
\subsection{13. Own estimations not done in the
paper}\label{own-estimations-not-done-in-the-paper}}

As an original extension, I estimated a TWFE-style interaction model
allowing the effect of the rate gap to differ for households above the
median balance-to-original-loan ratio. This is a simple way to test
whether rate lock is more visible among relatively high-balance
households.

\begin{verbatim}
                            OLS Regression Results                            
==============================================================================
Dep. Variable:              move_flag   R-squared:                       0.041
Model:                            OLS   Adj. R-squared:                  0.041
Method:                 Least Squares   F-statistic:                     364.8
Date:                Thu, 28 May 2026   Prob (F-statistic):               0.00
Time:                        17:15:49   Log-Likelihood:                -16506.
No. Observations:               48737   AIC:                         3.302e+04
Df Residuals:                   48731   BIC:                         3.308e+04
Df Model:                           5                                         
Covariance Type:                  HC1                                         
=========================================================================================
                            coef    std err          z      P>|z|      [0.025      0.975]
-----------------------------------------------------------------------------------------
const                    -0.0094      0.002     -5.166      0.000      -0.013      -0.006
rate_gap                 -0.0008      0.003     -0.313      0.754      -0.006       0.004
rate_gap_high_balance     0.1110      0.014      8.134      0.000       0.084       0.138
balance_to_orig       -1.436e-05   4.83e-06     -2.974      0.003   -2.38e-05    -4.9e-06
loan_age_years           -0.0531      0.001    -42.027      0.000      -0.056      -0.051
log_balance              -0.0041      0.001     -2.812      0.005      -0.007      -0.001
==============================================================================
Omnibus:                    14226.333   Durbin-Watson:                   2.037
Prob(Omnibus):                  0.000   Jarque-Bera (JB):            30955.793
Skew:                           1.760   Prob(JB):                         0.00
Kurtosis:                       4.690   Cond. No.                     1.77e+03
==============================================================================

Notes:
[1] Standard Errors are heteroscedasticity robust (HC1)
[2] The condition number is large, 1.77e+03. This might indicate that there are
strong multicollinearity or other numerical problems.
\end{verbatim}

The interaction term should be interpreted cautiously, but it offers a
clean example of an original panel-data extension beyond the baseline
replication workflow.

\hypertarget{list-of-references-cited-in-the-text}{%
\subsection{14. List of references cited in the
text}\label{list-of-references-cited-in-the-text}}

\begin{itemize}
\tightlist
\item
  Liebersohn, J., \& Rothstein, J. (2025). \emph{Household mobility and
  mortgage rate lock}. \emph{Journal of Financial Economics}.
  https://doi.org/10.1016/j.jfineco.2024.103973
\item
  Freddie Mac / Federal Reserve Bank of St.~Louis.
  \texttt{MORTGAGE30US}, 30-Year Fixed Rate Mortgage Average in the
  United States. https://fred.stlouisfed.org/series/MORTGAGE30US
\item
  Mendeley Data replication package:
  https://doi.org/10.17632/74sfv9kx9n.1
\item
  Replication Package README included locally in the project folder.
\end{itemize}

\hypertarget{appendix-to-do-list-for-the-next-panel-data-study}{%
\subsection{Appendix: To-do list for the next panel-data
study}\label{appendix-to-do-list-for-the-next-panel-data-study}}

\begin{itemize}
\tightlist
\item[$\square$]
  Define the economic question and the exact identifying variation
  before coding.
\item[$\square$]
  List the original paper's preferred dependent variable, key
  explanatory variable, controls, and fixed effects.
\item[$\square$]
  Verify raw data provenance and save permanent source links
  immediately.
\item[$\square$]
  Document every cleaning choice and every sample restriction in a
  reproducible script.
\item[$\square$]
  Check panel identifiers and time identifiers before any merge.
\item[$\square$]
  Count duplicates by individual-time cell before estimating anything.
\item[$\square$]
  Inspect missing values separately for dependent and explanatory
  variables.
\item[$\square$]
  Create the dynamic-panel benchmark sample with at least three
  consecutive observations.
\item[$\square$]
  Report attrition, holes, and sample-selection bias before regressions.
\item[$\square$]
  Compute between, within, FD and TWFE transformations early and compare
  their distributions.
\item[$\square$]
  Sort variables by within-variance share to understand which estimators
  are informative.
\item[$\square$]
  Check lags and first differences manually on the first stacked
  observations.
\item[$\square$]
  Inspect heterogeneity by individuals before relying on pooled
  coefficients.
\item[$\square$]
  Test robustness with alternative functional forms and subgroup
  interactions.
\item[$\square$]
  Assess instrument strength explicitly before reporting IV results.
\item[$\square$]
  Write interpretation in plain language after every table, not only
  code comments.
\item[$\square$]
  Keep an ordered folder of outputs: data, tables, figures, report,
  and work notes.
\item[$\square$]
  Export final tables in both machine-readable and human-readable
  formats.
\item[$\square$]
  State clearly what is replicated exactly and what is only
  approximated.
\item[$\square$]
  Finish with a limitations paragraph and a short agenda for future
  work.
\end{itemize}

\end{document}
